\section{Database Design}

\subsection{ER Diagram}
The database is designed around five core entities: Users, ExpenseCategories, BankAccounts, Income, and Expenses. The ER diagram models the relationships among users, their accounts, their income transactions, and their categorized expense transactions.

\subsection{Relational Schema}
The relational schema of the project is summarized as follows:
\begin{itemize}
    \item Users(UserID, UserName, Email, PhoneNumber)
    \item ExpenseCategories(CategoryID, CategoryName)
    \item BankAccounts(AccountID, UserID, BankName, Balance)
    \item Income(IncomeID, UserID, AccountID, Amount, IncomeDate, Description)
    \item Expenses(ExpenseID, UserID, CategoryID, AccountID, Amount, ExpenseDate, Description)
\end{itemize}

\subsection{Relationship Design}
A user can own multiple bank accounts, earn multiple income records, and make multiple expense records. Each expense belongs to one category, while each income and expense transaction is associated with one bank account.

\section{Database Implementation}

\subsection{Table Creation}
The database is implemented in MySQL. Primary keys are defined for each table, and foreign keys are added to preserve referential integrity between users, accounts, categories, income records, and expense records.

\subsection{Sample Data}
Representative sample data is inserted into every table to test constraints, joins, reports, and application functions. The sample data includes several users, multiple bank accounts, income transactions, expense categories, and expense transactions.

\subsection{Indexes}
Indexes are created on frequently queried columns such as UserID, AccountID, CategoryID, and transaction dates. These indexes improve the performance of financial summaries and expense lookup queries.

\subsection{Views}
Two views are created for reporting support:
\begin{itemize}
    \item MonthlySummaryView for monthly income, expense, and net amount,
    \item CategorySpendingView for category-wise expense analysis.
\end{itemize}

\subsection{Stored Procedures, Functions, and Triggers}
Stored procedures are used to insert income and expense transactions and to compute monthly summaries. User-defined functions are used to calculate total income, total expense, net amount, and budget status. Triggers automatically update bank account balances after new income or expense transactions are inserted.

\section{Python Application Development}

\subsection{Database Connection}
The Python application connects to MySQL by using the \texttt{mysql-connector-python} library. A separate connection module is prepared so the remaining program can reuse the same database configuration.

\subsection{Finance Tracker Functions}
The Python application supports adding income, adding expenses, viewing summaries, and checking budget status. These features interact directly with the MySQL database and reuse the database logic already implemented in SQL.

\subsection{Reporting}
The system can display financial summaries in tabular form and can also be extended with graphical reporting such as monthly bar charts or category-based charts.
